\documentclass[conference]{IEEEtran}

\usepackage[T1]{fontenc}
\usepackage{newtxtext}
\usepackage{cite}
\usepackage{amsmath,amssymb,amsfonts}
\usepackage{tabularx}
\usepackage{color}
\usepackage{tikz}
\usepackage{pgfplots}
\pgfplotsset{compat=1.18}
\usetikzlibrary{backgrounds,calc,positioning,arrows.meta,shapes.geometric,fit}
\usepackage{url}
\usepackage{textcomp}
\usepackage{algorithm}
\usepackage{algpseudocode}
\usepackage{float}
\usepackage[urlcolor=blue]{xurl}
\sloppy

\newcommand{\todo}[1]{\textcolor{red}{\textbf{[TODO: #1]}}}

\definecolor{hS}{RGB}{134,199,134}
\definecolor{hP}{RGB}{245,160,110}
\definecolor{hW}{RGB}{233,128,118}
\definecolor{hN}{RGB}{236,236,236}

\newcolumntype{L}[1]{>{\raggedright\arraybackslash}p{#1}}
\def\lblS{\textbf{S}}
\def\lblP{\textbf{P}}
\def\lblW{\textbf{W}}
\def\lblN{\textbf{-}}

\newcommand{\hcw}{1.25}
\newcommand{\hch}{0.50}

\newcommand{\heatrow}[3]{
  \node[anchor=east, font=\small] at (4.5,{-(#1-0.5)*\hch}){#2};
  \foreach \r [count=\j] in {#3}{
    \node[draw=white,line width=0.8pt,fill=h\r,minimum width=\hcw cm,minimum height=\hch cm,inner sep=0pt]
    at ({4.65+(\j-0.5)*\hcw},{-(#1-0.5)*\hch}) {\csname lbl\r\endcsname};
  }
}

\begin{document}

\title{Trust Boundaries and Blind Spots: A Taxonomy of AI Supply Chain Attack Surfaces and Defensive Readiness Assessment}

\author{
\IEEEauthorblockN{Shanu Shekhar Giri}
\IEEEauthorblockA{\textit{Cyber Security Integrated Specialization M.Tech.} \\
VIT Bhopal University\\
Sehore, Madhya Pradesh, India \\
shanushekhargiri2022@vitbhopal.ac.in}
\and
\IEEEauthorblockN{Dr. Adarsh Patel}
\IEEEauthorblockA{\textit{Senior Assistant Professor, CSDF, SCAI} \\
\textit{VIT Bhopal University}\\
Sehore, Madhya Pradesh, India \\
Adarsh.patel@vitbhopal.ac.in}
\and
\IEEEauthorblockN{Dr. Sajjad Ahmed}
\IEEEauthorblockA{\textit{Assistant Professor, CSDF, SCAI} \\
\textit{VIT Bhopal University}\\
Sehore, Madhya Pradesh, India \\
sajjad.ahmed@vitbhopal.ac.in}
\and
\IEEEauthorblockN{Dr. Praveen Lalwani}
\IEEEauthorblockA{\textit{Assistant Professor, CSDF, SCAI} \\
\textit{VIT Bhopal University}\\
Sehore, Madhya Pradesh, India \\
praveen.lalwani@vitbhopal.ac.in}
}

\maketitle
\raggedbottom

\begin{abstract}
Third-party artifacts---pretrained models, web-scale datasets, plugins, agent skills, and software libraries---are the building blocks of AI systems and their supply chains. We identify three ways this risk differs from traditional software supply chains: models are executable artifacts; data shapes behavior without being code; and natural-language instructions act as executable directives. We propose a four-layer taxonomy of AI supply chain attack surfaces, analyze five documented incidents, and rate eight defenses against thirteen attack classes. Across the incidents we find three failure modes: stale checks (time-of-check/time-of-use), interpretation differentials (a scanner and a loader parse a file differently), and absent or blind controls. To examine the second mode, we build \texttt{pickle-harness}, a network-isolated differential-testing harness that compares scanner verdicts with observed execution. On a synthetic 7,500-file corpus, the standard loader executed all 5,000 malicious files. ModelScan (v0.8.4) flagged all of them. Picklescan (v0.0.22) flagged 3,750 and ended in a parse error on the 1,250 nullifAI-style files; the output for every one of them also reports the payload's dangerous global, so whether this counts as a detection depends on how a pipeline interprets a mixed verdict. None of the files executed under \texttt{torch.load(weights\_only=True)} or \texttt{safetensors}, but because they were not valid containers for those loaders, this measures format rejection and not the strength of the restricted unpickler. The coverage ratings are analytical estimates, and the claim that removing the execution mechanism outperforms inspecting for it remains a hypothesis that our experiments support only partially.
\end{abstract}

\begin{IEEEkeywords}
AI supply chain security, malicious models, data poisoning, model context protocol, tool poisoning, empirical evaluation, pickle harness
\end{IEEEkeywords}

\section{Introduction}

Most artificial intelligence systems today are not built from scratch. Developers are able to obtain pretrained weights from public hubs. They train on third-party scraped data. They provide language-model agents with external tools and plugins. They install the libraries that all this is based on.

The AI supply chain is increasingly in the crosshairs. Every dependency gives an attacker a point of injection for malicious material before the system goes into production. The reported cases are troubling on all four levels of reliance.

In February 2024, JFrog found about 100 malicious models on Hugging Face~\cite{jfrog2024}. The majority were reverse-shell payloads in PyTorch's pickle format. In 2025, ``nullifAI'' models bypassed Hugging Face's Picklescan, using deliberately corrupted pickle streams~\cite{reversinglabs2025}. Carlini \textit{et al.} showed that poisoning 0.01\% of LAION-400M or COYO-700M would cost around \$60 at the dataset level~\cite{carlini2024}. In April 2025, Invariant Labs showed how a malicious Model Context Protocol (MCP) server could use instructions embedded in a tool description to steal a user's WhatsApp history~\cite{invariant2025}. In September 2025, the \texttt{postmark-mcp} package was found to be secretly copying users' outgoing email to an attacker~\cite{koi2025}. In February 2026, researchers found 341 malicious skills in the ClawHub agent registry~\cite{csa2026}. CVE-2025-6514 was a bug at the dependency layer in the npm package \texttt{mcp-remote} that allowed a malicious MCP server to execute commands on the client host~\cite{nvd2025}.

These cases have mostly been studied in isolation. They include scanner evasion on model files, dataset poisoning, and MCP attack vectors. More extensive work on the large language model (LLM) supply chain, along with industry groups such as the Cloud Security Alliance (CSA), have started to categorize threats at the agent layer. However, to the best of the authors' knowledge, none of them evaluates defensive coverage against reported incidents across all layers. The domain is also faster than the peer-reviewed literature. So this study draws on academic work as well as vendor disclosures. Where the quality of evidence matters, it tells you which is which.

\begin{figure}[!t]
\centering
\begin{tikzpicture}[transform shape,scale=0.55]
\draw[thick, fill=blue!20, fill opacity=0.6] (0,0.8) circle (1.5);
\draw[thick,fill=orange!30, fill opacity=0.6] (-1.0,-0.8) circle (1.5);
\draw[thick,fill=green!20, fill opacity=0.6] (1.0,-0.8) circle (1.5);
\node[align=center,font=\bfseries] at (0,2.6) {Academic Research};
\node[align=center,font=\bfseries] at (-2.4,-1.6) {Disclosures from Vendors};
\node[align=center,font=\bfseries] at (2.4,-1.6) {Industry\\Standards};
\node[star,star points=5,star point ratio=2.25,fill=yellow,draw=black,inner sep=1.5pt] (star) at (0,-0.2) {};
\node[draw, rounded corners=3pt, fill=yellow!20, inner sep=4pt, text width=5cm, align=center, font=\small] (gap) at (0,-4.2) {\textbf{Gap:} unified cross-layer taxonomy and defense coverage assessment};
\draw[-{Stealth[length=2.5mm]}, thick] (gap.north) -- (star.south);
\end{tikzpicture}
\caption{AI supply chain fragmentation. Academic studies, vendor disclosures, and industry recommendations each tackle a different piece of the problem. The star points to the unified gap considered in this paper.}
\label{fig:venn}
\end{figure}

We analyze the major AI system supply chain attack surfaces and the extent to which existing defenses cover them, as illustrated in Fig.~\ref{fig:venn}. Our contributions are fourfold.

The first is a taxonomy. It is a four-level taxonomy of AI supply chain attack surfaces, ordered by compromised artifact.

The second is incident analysis. Five reported incidents were studied to determine which control had failed and why. This leads to three failure modes.

Third is coverage assessment. It is an analytical rating of 13 attack classes against 8 defenses. It distinguishes mechanism elimination, runtime verification, and point-in-time inspection.

The fourth is an empirical case study. We compare two static pickle scanners against observed execution using \texttt{pickle-harness}. We use it to test scanner fail-open behavior, to compare to a baseline (Song \textit{et al.}) and other previous reports, and to characterize what a format-rejection result demonstrates and does not demonstrate.

The remainder of the paper is organized as follows. Section~\ref{sec:background} describes background and related work. Section~\ref{sec:taxonomy} introduces the taxonomy. Incidents are reported in Section~\ref{sec:cases}. Section~\ref{sec:case-observations} describes the cross-case observations. The defenses are set out in Section~\ref{sec:defenses}. The description of the harness is in Section~\ref{sec:method}. Its results are given in Section~\ref{sec:eval}. Open problems and limitations are in Section~\ref{sec:discuss}, and the paper concludes in Section~\ref{sec:conclusion}.

\section{Background and Related Work}
\label{sec:background}

\subsection{Attacks on the software supply chain}
\label{sec:sw-supply-chain}

A software supply chain attack infects a dependency before delivery to the target, so trusted software ships attacker code.

Typosquatting is publishing a malicious package with a name similar to a popular one, like \texttt{reqeusts} for \texttt{requests}. It is caused by a typo in the install command. With dependency confusion, the attacker publishes a public package with the same name as the private package, but with a higher version number. The resolver then prefers the public artifact. Both are based on running code at install time, where fetching a package is already an execution event. This is the baseline threat model that the AI pipeline builds on and extends.

\subsection{The AI Pipeline}
\label{sec:ai-pipeline}

An AI system is constructed by a pipeline of multiple steps, illustrated in Fig.~\ref{fig:pipeline}. Each step introduces a third-party dependency. Training data is third-party. Training creates a serialized model artifact. That artifact is available through a registry. Then the deployed model is augmented with agents and plugins that call external tools at runtime. The pipeline is thus a chain of trust boundaries. A compromise at one boundary carries into the stages beyond it. Protection that works for one type of artifact may not work for the next.

\begin{figure*}[!t]
\centering
\begin{tikzpicture}[
font=\footnotesize,
stage/.style={rectangle,rounded corners=4pt,draw=black!70, thick, minimum width=2.8cm, minimum height=1.2cm, align=center, fill=white},
surf/.style={rectangle,rounded corners=3pt,draw=red!70!black, thick, fill=red!10, minimum width=3.8cm, minimum height=0.8cm, align=center},
arr/.style={-{Stealth[length=3mm, width=2mm]}, thick, black!80}
]
\node[stage,fill=blue!10] (d) at (0,0) {\textbf{Data acquisition}};
\node[stage, fill=green!10] (t) at (3.5,0) {\textbf{Fine-tuning}\\training};
\node[stage, fill=orange!10] (m) at (7,0) {\textbf{Model}\\serialized artifact};
\node[stage, fill=violet!10] (dep) at (10.5,0) {\textbf{Deployment}\\serving};
\node[stage, fill=teal!10] (a) at (14,0) {\textbf{Agents/plugins}\\tools, MCP};

\foreach \p/\q in {d/t,t/m,m/dep,dep/a} \draw[arr] (\p.east) -- (\q.west);

\node[surf] (s2) at (0,2.2) {\textbf{Datasets}\\split view, RAG};
\node[surf] (s1) at (7,2.2) {\textbf{Models}\\backdoors, pickle};
\node[surf] (s3) at (14,2.2) {\textbf{Plugins \& tools}\\poisoning, rug pulls};

\draw[red!70!black, thick, dashed] (s2.south) -- (d.north);
\draw[red!70!black, thick, dashed] (s1.south) -- (m.north);
\draw[red!70!black, thick, dashed] (s3.south) -- (a.north);

\node[font=\scriptsize\itshape, text=black!70] at (0,-1.2) {Data shapes behavior};
\node[font=\scriptsize\itshape, text=black!70] at (7,-1.2) {Models are executable};
\node[font=\scriptsize\itshape, text=black!70] at (14,-1.2) {Language is code};

\node[rectangle,draw=gray!80, fill=gray!15, rounded corners=3pt, minimum width=16.5cm, minimum height=0.8cm, align=center] (deps) at (7,-2.5) {\textbf{Dependencies:} Packages, ML frameworks, protocol SDKs (transitive trust across all stages)};
\end{tikzpicture}
\caption{The AI pipeline and its four attack surfaces. Artifact-specific surfaces are shown in red boxes. The gray bar shows the dependency layer that exists in all stages. Italic labels indicate the distinguishing property of each artifact.}
\label{fig:pipeline}
\end{figure*}

\subsection{What Makes AI Different}
\label{sec:why-ai-differs}

Artificial intelligence systems are exposed to the ordinary supply chain risks described above, but three properties make them harder to defend.

\emph{Models are executable artifacts.} PyTorch checkpoints have historically been stored in the Python pickle format. This format serializes objects into a stream of opcodes that are interpreted by the loader. The opcodes can import and invoke arbitrary callables. Thus loading an untrusted pickle is as dangerous as executing an untrusted program. The format is still in use. An analysis of 10,023 popular Hugging Face repositories created between July 2025 and June 2026 showed that 9.3\% depend on pickle~\cite{qin2026}. The inherent format hazard makes it difficult to make inspection-based defenses complete. Liu \textit{et al.} identified 22 vulnerable pickle load paths, and nineteen of them got past all four scanners they tested~\cite{liu2025}.

\emph{Data defines behavior, but it is not code.} A poisoned dataset accomplishes nothing on its own. It alters what a model learns. Carlini \textit{et al.} demonstrated that small and inexpensive injections into web-scale datasets are possible. This is a provenance and integrity problem, not a code-execution problem. The signing infrastructure built for models has not been extended to datasets.

\emph{Code is natural-language instructions.} An agent converts descriptions of third-party tools and tool replies into part of its instruction context. Tool descriptions are considered only once, when an agent connects. Tool responses enter the model context without equivalent validation. Empirical work on the MCP ecosystem shows that malicious servers can manipulate agents via tool metadata and returned content, with no vulnerable code~\cite{song2025}. Any instruction in either channel is a command, and the trust persists once a server is approved.

\subsection{Review of Literature}
\label{sec:literature}

Research on AI supply chain security is scattered over a number of communities. Most of the studies deal with one layer of the pipeline. We survey work on plugins and protocols, model files, datasets, cross-layer surveys, and provenance and signing, and end with a comparison table.

\textbf{Plugin and protocol layer.} One of the most systematic studies of the MCP attack surface is presented by Song \textit{et al.}~\cite{song2025}. They classify attacks into four categories: tool poisoning, puppet attacks, rug pulls, and exploitation of malicious external resources. They show them on MCP aggregation platforms. In a user study, they discovered that users often did not recognize malicious servers. They present a taxonomy of attacks and results on their feasibility, but not an evaluation of defenses.

\textbf{Model layer.} PickleBall~\cite{pickleball2025} statically analyzes the source code of machine learning libraries. It computes a per-model loading policy, and enforces it at load time. It loaded 79.8\% of the 252 benign models, while rejecting all 84 malicious ones. It loaded 22\% more benign models than the state-of-the-art loader it was compared to. The same evaluation found that the scanners tested did not detect known malicious models. Qin \textit{et al.}~\cite{qin2026} reported high false-positive rates for current scanners. They propose a context-aware alternative, but the results are self-reported on their own benchmark. PickleBall narrows the execution path rather than eliminating it.

\textbf{Dataset layer.} Plant \textit{et al.}~\cite{plant2025} investigate poisoning of BERT in low-resource settings. The first such study for Basque is included in this. They observe average performance drops above 30\% with poisoning ratios greater than 50\%. Carlini \textit{et al.}~\cite{carlini2024} ask a different question: how cheap can poisoning web-scale data be in practice? They show that \$60 would have been sufficient to poison 0.01\% of LAION-400M or COYO-700M using split-view and frontrunning attacks. The two studies together establish the effect and feasibility of poisoning.

\textbf{Cross-layer surveys.} De Maio \textit{et al.}~\cite{demaio2025} survey LLM vulnerabilities across the training, inference and deployment lifecycle. They analyze 17 vulnerabilities, such as dependency risks in retrieval-augmented generation systems. Deployment is the least studied phase, they note, with just 10 of the 66 selected papers dealing with it. They call for provenance tracking and verifiability in LLM pipelines. Berini \textit{et al.}~\cite{berini2026} present an end-to-end taxonomy of data-centric, model-centric, and deployment-centric threats. Their taxonomy explicitly addresses supply chain vulnerabilities and plugin interfaces.

\textbf{Provenance and signing.} Campbell~\cite{campbell2026} proposes a Model Bill of Materials with Post-Quantum-Safe Extensions (MBOM-PQC). The proposal includes a single signing and attestation pipeline. It also contains a five-level Supply Chain Assurance Maturity Model (SCAMM). This closes the provenance gap by specifying a verifiable schema for AI artifacts. However, it is a design-science contribution, which is offered for future operational validation. It has not been assessed across heterogeneous model ecosystems, and there is no equivalent mechanism for datasets.

\subsection{Placement and Comparison}
\label{sec:positioning}

Table~\ref{tab:litcompare} summarizes what each reviewed paper produced and how this paper builds on it. To our knowledge, no previous study has rated a common set of defenses against classes of attack in all four layers, and no previous research connects each rating to reported incidents. This paper provides that rating as an analytical judgment. It adds a small empirical harness (\texttt{pickle-harness}) that reproduces a scanner/loader mismatch in a controlled setting for the model layer. We take Song \textit{et al.} as a baseline for the plugin layer.

\begin{table*}[!t]
\caption{What each paper produced versus what this paper contributes}
\label{tab:litcompare}
\centering
\footnotesize
\renewcommand{\arraystretch}{1.25}
\begin{tabular}{L{2.2cm}L{2.0cm}L{3.3cm}L{5.5cm}}
\hline\hline
\textbf{Paper} & \textbf{Layer} & \textbf{What they produced} & \textbf{How this paper extends it}\\
\hline
Song \textit{et al.}~\cite{song2025} & Plugins / MCP & Four-class MCP attack taxonomy; feasibility on aggregation platforms; user study on server recognition & Takes their attack classes as input and rates defensive coverage for all four, linked to the Invariant WhatsApp incident\\
PickleBall~\cite{pickleball2025} & Models & Policy-based secure pickle loader; 79.8\% benign load rate; 100\% malicious rejection & Replicates a mixed scanner verdict on nullifAI-style files using \texttt{pickle-harness}; ranks mechanism elimination (safetensors) as structurally stronger than policy enforcement\\
Plant \textit{et al.}~\cite{plant2025} & Datasets & BERT performance degradation (over 30\% at over 50\% poisoning) in low-resource settings & Uses it as impact evidence for the dataset layer; notes that provenance defenses are largely absent for datasets\\
Carlini \textit{et al.}~\cite{carlini2024} & Datasets & Split-view and frontrunning attacks; 0.01\% poisoning for \$60 & Uses their feasibility and cost results as input; rates curation-time hashing and snapshot pinning as partial or weak in the coverage matrix\\
De Maio \textit{et al.}~\cite{demaio2025} & Cross-layer & Lifecycle survey of LLM vulnerabilities (17); deployment phase under-studied & Gives an artifact-centric, defense-centric coverage matrix over four layers; the deployment-phase gap motivates the plugin-layer scope\\
Berini \textit{et al.}~\cite{berini2026} & Cross-layer & Lifecycle threat taxonomy; maps threats to defenses; covers supply chain flaws & Moves from threat-to-defense mapping to an ordinal coverage rating per attack class, based on incidents, with one empirical scanner experiment\\
Campbell~\cite{campbell2026} & Provenance / Dependencies & MBOM-PQC schema; PQC-safe signing pipeline; SCAMM maturity model & Rates signing as partial against malicious files and weak against backdoors, as a verified publisher can still ship malicious weights\\
\hline
\end{tabular}
\end{table*}

\subsection{Important Concepts}
\label{sec:key-concepts}

\emph{Safetensors.} Safetensors saves the raw tensor data behind a JSON header. It has no execution path at load time. It is meant to replace pickle for model distribution. It does not cover other formats that can embed executable code, such as Keras Lambda layers, which Hugging Face separately inspects with scanners.

\emph{Model Context Protocol.} MCP is a protocol for how LLM applications can talk to external tools and data sources using JSON-RPC. It works between hosts, clients and servers. Servers expose tools, resources and prompts. The specification assumes tool output is not to be trusted, but clients usually put it in the model context without validation. This channel is exploited by tool poisoning.

\emph{Agent skills.} Skills are bundles of versioned text organized around an \texttt{SKILL.md} file. They are published through registries like ClawHub. A skill can be documentation only or may include executable code. A registry-distributed skill is a supply chain artifact, similar to a package.

\section{AI Supply Chain Attack Surface Taxonomy}
\label{sec:taxonomy}

We organize attack surfaces by the artifact an attacker compromises, summarized in Table~\ref{tab:taxonomy}. Each layer is tied to a property from Section~\ref{sec:why-ai-differs}. Models are tied to executability. Datasets are tied to behavioral effect. Plugins and tools are tied to instruction-as-code. Dependencies are tied to inherited trust.

\begin{table*}[!t]
\caption{Taxonomy of AI supply chain attack surfaces}
\label{tab:taxonomy}
\centering
\footnotesize
\renewcommand{\arraystretch}{1.2}
\begin{tabular}{L{2.0cm}L{3.4cm}L{3.1cm}L{4.0cm}L{3.5cm}}
\hline
\textbf{Layer} & \textbf{Compromised artifact} & \textbf{Exploited AI property} & \textbf{Attack classes} & \textbf{Anchor incident}\\
\hline
Models & Serialized model files and weights & Models are executable artifacts & Malicious pickle; backdoored weights & JFrog (2024); nullifAI (2025)\\
Datasets & Training corpora, RAG knowledge bases & Data shapes behavior without being code & Split-view, frontrunning, and RAG poisoning & Carlini \textit{et al.} (2024)\\
Plugins and tools & MCP servers, tool definitions, agent skills & Language acts as code & Tool poisoning; rug pulls; tool squatting; malicious skills & Invariant WhatsApp (2025); ClawHub (2026)\\
Dependencies & Packages, frameworks, protocol SDKs & Inherited transitive trust & Typosquatting; dependency confusion; unsafe tooling & postmark-mcp (2025); CVE-2025-6514\\
\hline
\end{tabular}
\end{table*}

\subsection{Models: Malicious Serialized Files and Backdoored Weights}
\label{sec:tax-models}

Model artifacts can be compromised in two ways. \emph{Malicious serialization} places code in a pickle stream that executes at load time; JFrog found roughly 100 such models in 2024. NullifAI showed that a scanner and a loader can disagree about the same file when the stream is altered. This mismatch is systematic, because scanners may miss the pickle stream that the framework will actually load.

\emph{Backdoored weights} load correctly but behave badly on a trigger. They arise from poisoned training data or from fine-tuning. They are harder to detect than serialization payloads, because no code runs at load time and the weights are valid tensors. The two techniques are complementary: serialization gives immediate code execution, while a backdoor gives long-term control over model behavior.

\subsection{Datasets: Frontrunning, Split-View, Retrieval Poisoning}
\label{sec:tax-datasets}

Carlini \textit{et al.} define two dataset attacks that are feasible today. \emph{Split-view poisoning} exploits the mutability of online content: the data an annotator saw differs from the data a later client downloads, so an attacker who controls a source site can change it after curation. \emph{Frontrunning poisoning} targets datasets that are periodically snapshotted from crowd-sourced content such as Wikipedia, and needs only a short window between an edit and the snapshot.

Retrieval-augmented generation (RAG) moves the dataset surface from training time to inference time. Poisoning the knowledge base requires no access to training, because the knowledge base is read at query time. Poisoning attacks have been shown to succeed across retrievers, including black-box systems, and against several defenses. We treat RAG poisoning as a sub-category of datasets because the compromised artifact is still a corpus. Unlike pre-training, where the attack window is bounded, the window here is continuous.

\subsection{Tools and Plugins: Malicious Skills, Squatting, Rug Pulls and MCP Poisoning}
\label{sec:tax-plugins}

Song \textit{et al.}~\cite{song2025} study MCP attacks end to end and classify them into four types: \emph{tool poisoning}, \emph{puppet attacks}, \emph{rug pulls}, and \emph{exploitation via malicious external resources}. They implement malicious servers on MCP aggregation platforms, and show in a user study that users had difficulty spotting them.

In tool poisoning, an adversarial instruction is placed in a tool description or other metadata. The agent reads it as natural language and trusts it as guidance for its reasoning. \emph{Rug pulls} add a temporal dimension: a server publishes benign definitions, obtains approval, and then changes them. In a \emph{puppet attack}, a malicious server's crafted tool description steers the agent's decisions about which tools to call, in an environment where several servers coexist, and the malicious server need not be invoked itself. In \emph{exploitation via malicious external resources}, the server's code and descriptions are benign, but a tool fetches third-party content, such as a web page, that carries hidden instructions. \emph{Tool squatting} exploits the lack of a namespaced tool registry: a malicious server registers a tool name identical to one on a trusted server, and an agent that selects tools by name may run the wrong implementation.

Agent skill registries present a parallel surface. A malicious skill is a supply chain artifact like a package, and the discovery of 341 malicious skills in ClawHub indicates that registry moderation is immature at this layer.

\subsection{Dependencies: Typosquatted or Compromised Packages and Vulnerable Tools}
\label{sec:tax-dependencies}

Typosquatting, dependency confusion, and compromised transitive libraries apply to ML frameworks, tokenizers, and agent tools. The dependency surface of Section~\ref{sec:sw-supply-chain} therefore carries over to AI systems. What is new is that tool-layer protocol implementations now execute logic as a design convenience. For example, in CVE-2025-6514 a malicious MCP server could supply a crafted \texttt{authorization\_endpoint} URL to \texttt{mcp-remote} that ran commands on the client host. Vulnerable tooling is therefore not a separate category from dependency compromise. It is the form dependency compromise takes when the dependent is an agent protocol implementation.

\section{Case Studies}
\label{sec:cases}

Five incidents are analyzed, covering all four layers. The first is malicious pickle models (models). The second is web-scale dataset poisoning (datasets). Third is MCP tool poisoning (plugins and tools). The fourth is postmark-mcp (dependencies). Fifth is ClawHub (agent skills, a registry-distributed artifact). Each case is described by what happened, the mechanism, the impact and detection, and what went wrong. There are only a few in-the-wild incidents. The Invariant demonstration and the Carlini \textit{et al.} attacks are research proofs of concept.

\subsection{Malicious Pickle Models: JFrog and nullifAI}
\label{sec:casemodels}

In February 2024, JFrog found some 100 malicious models on Hugging Face. Some 95\% were in PyTorch's pickle format, with reverse-shell payloads that fire when the model is loaded. In early 2025, ReversingLabs found two other models that bypassed Hugging Face's Picklescan using what it described as the nullifAI technique.

In both cases the malicious logic was embedded in the pickle \texttt{\_\_reduce\_\_} mechanism. The nullifAI models were compressed with 7z instead of the normal ZIP. They purposely had a broken stream, and they put the payload at the front. The payload ran before the stream failed, because pickle deserializes sequentially, and the scanner did not catch the file.

The code runs with the same privileges as the process loading the model, which is typically a data scientist's workstation or a build runner. Both sets of models were discovered by outside researchers. Hugging Face took down the nullifAI models and improved Picklescan after the disclosure.

The file is parsed differently by the loader and the scanner. A stream can run its payload before a parser decides it is not acceptable. The upload-time scanner checks its own parse of the file, not what the loader will execute. Safetensors removes the execution path, but pickle-only repositories are still common.

\subsection{Poisoning Web-Scale Datasets: LAION, COYO and Wikipedia}
\label{sec:case-datasets}

Carlini \textit{et al.} demonstrated two practical poisoning attacks. One targeted datasets used to train production models (LAION-400M, COYO-700M). The other targeted snapshot datasets built from crowd-sourced content such as Wikipedia.

In \emph{split-view poisoning}, a source URL controlled by the attacker is modified after curation, so the annotation and the downloaded content are out of sync. In a \emph{frontrunning} poison, the attacker manipulates the crowd-sourced content just before the snapshot is taken, and the poisoned version becomes the canonical training example.

The authors state that poisoning 0.01\% of LAION-400M or COYO-700M would cost about \$60. The model trains normally, and the effect is only in the trained model. The authors informed the parties involved and proposed low-overhead defenses. There was no runtime detection event.

There was no integrity check in the curation. The datasets were released as lists of URLs without content hashes or signed manifests. Nothing distinguished the downloaded content from the annotated content. The gap is a provenance record that is missing, not a scanner that is missing.

\subsection{MCP Tool Poisoning: WhatsApp Demonstration by Invariant Labs}
\label{sec:case-mcp-tool}

In April 2025, Invariant Labs showed how a malicious MCP server could exfiltrate the user's chat history. The malicious server was associated with the same agent as a genuine WhatsApp server. To do this it manipulated the description of its tool.

The attack was a combination of rug pull and tool shadowing. The malicious server's description looked harmless at first, and the user's approval was obtained. The description later instructed the agent to send \texttt{send\_message} calls to the separate and legitimate \texttt{whatsapp-mcp} server. The calls were directed to a number controlled by the attacker, with the chat history in the message body. The payload had been padded with whitespace, so the confirmation dialog only displayed a short, innocuous message.

The chat history left through the normal WhatsApp channel without the malicious tool being directly called. The whitespace padding defeated the only control, the confirmation dialog. The flaw was discovered and published by the researchers themselves. Approval was a one-time check of the description at connection, and the revised description was not checked against the approved version.

\subsection{postmark-mcp: A Backdoored MCP Package}
\label{sec:case-postmark}

In September 2025, researchers observed that the npm package \texttt{postmark-mcp} had been modified in version 1.0.16. The package looked like a normal Postmark integration. The modification redirected all outgoing emails to an attacker-controlled mailbox.

For 15 versions the package acted like a working MCP server and gained downloads and trust. One of the releases added the attacker's address as a hidden recipient. It took the code from the real Postmark repository and published it under a very similar name. The attack was impersonation and a malicious upgrade, not a protocol exploit.

The backdoor could quietly copy password resets, invoices and internal documents. The package was downloaded over 1,600 times before being pulled. The discovery came from a third-party monitoring service, not from the package registry or its users.

Reputation and registry vetting failed, because one bad version inherited the trust of 15 good ones. There was no automated check to see whether the new version behaved the same as the previous versions, particularly for its outbound data flows. A consumer who had relied on an earlier, audited version would have been exposed to the change.

\subsection{ClawHub: Malicious Agent Skills}
\label{sec:case-clawhub}

An audit of the ClawHub registry in February 2026 found 341 malicious skills out of 2,857, or around 12\%. 335 of them were part of a single campaign, dubbed ClawHavoc.

The skills were advertised as tools for crypto traders, content creators and financial professionals. In the ``prerequisites'' section, they instructed users to download and run a password-protected archive (Windows) or an encoded script (macOS). The chain served Atomic Stealer on macOS, which steals passwords, browser data, crypto wallets and SSH keys. The attack did not take advantage of a vulnerability, but used the normal skill installation workflow.

The targets were high-value individuals running always-on agents with access to wallets and credentials. OpenClaw later introduced community reporting that automatically hides a skill after three reports.

Registry admission control failed. The malicious step was a user-run prerequisite, not skill code. Uploads were accepted without behavioral review. Static analysis cannot differentiate between a skill that calls an API and one that also exfiltrates credentials.

\section{Cross-Case Observations}
\label{sec:case-observations}

The control that failed in each case is listed in Table~\ref{tab:incidents}. There are three modes of failure.

\emph{Stale checks (time-of-check/time-of-use).} Postmark-mcp and Invariant WhatsApp succeeded because a check was performed at time $t_1$ and the artifact changed at time $t_2$. The check never ran again.

\emph{Interpretation differentials (parser mismatch).} NullifAI involved no change after the check. The scanner and the loader read the same bytes at the same time and came to different conclusions (Fig.~\ref{fig:differential}). The loader ran the payload, but the scanner choked on the structure of the file.

\emph{Absent or blind controls.} LAION/COYO had zero integrity check. The static review of ClawHub could not see the prerequisite run by the user. In both cases the control was either non-existent or could not look at the relevant action.

We separate these modes because the remedies are different. Stale checks need to be re-verified at the point of use, either by re-validation or by hashing. Differentials have to clear up the ambiguity: either they remove the execution path, or they enforce a strict format that both the scanner and the loader understand. Missing or blind controls require provenance mechanisms or enhanced inspection. The three-way split is based on five events selected to cover all four levels, and is not meant to be exhaustive.

\begin{figure}[!t]
\centering
\begin{tikzpicture}[font=\footnotesize,
box/.style={rectangle,rounded corners=3pt,draw=black!70,thick,align=center,minimum height=0.9cm,minimum width=2.4cm},
arr/.style={thick,-{Stealth[length=2mm]}}]
\node[box,fill=orange!15] (f) at (0,0) {Same file\\(same bytes)};
\node[box,fill=blue!10] (s) at (4.2,0.9) {Scanner parse:\\stops at corruption};
\node[box,fill=red!15] (l) at (4.2,-0.9) {Loader parse:\\payload executes};
\draw[arr] (f.east) -- (s.west);
\draw[arr] (f.east) -- (l.west);
\end{tikzpicture}
\caption{Interpretation differential. The scanner and the loader are reading the same file at the same time, but getting different results. The attack is successful without any change after approval.}
\label{fig:differential}
\end{figure}

\begin{table*}[!t]
\caption{Failure of control in each case examined}
\label{tab:incidents}
\centering
\footnotesize
\renewcommand{\arraystretch}{1.2}
\begin{tabular}{L{1.8cm}L{3.0cm}L{3.4cm}L{7.2cm}}
\hline\hline
\textbf{Layer} & \textbf{Incident} & \textbf{Failed control} & \textbf{Why it failed}\\
\hline
Models & JFrog / nullifAI & Upload-time scanning & Scanner and loader read the file differently; the scanner's parser rejects the stream after the payload has already executed. (Differential)\\
Datasets & LAION/COYO poisoning & Integrity at curation time & No content hash or signed manifest, so later changes to the referenced content could not be detected. (Absent or blind)\\
Plugins & Invariant WhatsApp & Connect-time approval & Tool description modified after approval without re-validation; payload obfuscated by whitespace padding. (Stale)\\
Packages & postmark-mcp & Reputation and registry vetting & 15 clean releases preceded one malicious update; no comparison of version behavior. (Stale)\\
Skills & ClawHub & Registry admission control & Open upload; static analysis could not capture the user-run prerequisite that performed the damaging step. (Absent or blind)\\
\hline
\end{tabular}
\end{table*}

\section{Coverage and Defenses}
\label{sec:defenses}

\subsection{Evaluation Framework}
\label{sec:def-frame}

We categorize defenses according to time. A \emph{point-in-time check} examines an artifact once, such as a scan at upload, an approval at connection, or reputation built up over previous versions. \emph{Runtime verification} verifies the artifact during use. \emph{Mechanism elimination} takes away a capability the attacker needs, so no check needs to be done.

A point-in-time check passes at $t_1$ and has no knowledge that anything has changed at $t_2$. At $t_3$, a runtime check compares the artifact with the approved state, as illustrated in Fig.~\ref{fig:threat}.

\begin{figure*}[!t]
\centering
\begin{tikzpicture}[
font=\small,
box/.style={rectangle,rounded corners=3pt,draw=black!70, thick, inner sep=4pt, align=center,minimum height=1cm},
arr/.style={thick,-{Stealth[length=2.5mm, width=2mm]}}
]
\node[box, fill=blue!10, minimum width=5cm] (t0) at (0,0) {\textbf{\(t_0\):} Benign artifact published};
\node[box,fill=green!10, minimum width=5cm] (t1) at (0,-1.8) {\textbf{\(t_1\):} Approved, scanned and reviewed};
\node[box, fill=red!15, minimum width=5cm] (t2) at (0,-3.6) {\textbf{\(t_2\):} Artifact quietly edited (rug pull)};
\node[box, fill=orange!15, minimum width=5cm] (t3) at (0,-5.4) {\textbf{\(t_3\):} Agent loads modified artifact};
\foreach \a/\b in {t0/t1,t1/t2,t2/t3} \draw[arr] (\a.south) -- (\b.north);
\node[box, fill=red!10, minimum width=2.5cm, minimum height=1.5cm] (att) at (-5,-1.8) {\textbf{Attacker}\\ publishes, then\\ waits};
\draw[arr,red!70!black] (att.east)--(t0.west) node[font=\scriptsize,above,pos=0.5] {publish};
\draw[arr, red!70!black] (att.east) -- (t2.west) node[font=\scriptsize,above,pos=0.5] {change};
\node[box, fill=blue!10, minimum width=3.5cm] (d1) at (5.5,-1.8) {\textbf{Point-in-time check}\\scanner, one-time approval};
\node[box, fill=green!15, minimum width=3.5cm] (d2) at (5.5,-5.4) {\textbf{Runtime verification}\\hash re-check, gateway};
\draw[arr,blue!70!black] (d1.west) -- (t1.east) node[pos=0.5, above, font=\tiny] {passes};
\draw[arr,dashed, red!70!black] (d1.south) -- (t2.east) node[pos=0.5,right,font=\scriptsize,align=left] {not\\re-run};
\draw[arr,green!60!black] (d2.west) -- (t3.east) node[pos=0.5, above, font=\scriptsize] {mismatch: block};
\end{tikzpicture}
\caption{Threat model for time-delayed attacks (the stale-check mode). The check at time $t_1$ passes, but the change at $t_2$ is missed because the check is not repeated. At $t_3$ the artifact is blocked because of a hash mismatch against the approved hash.}
\label{fig:threat}
\end{figure*}

Coverage is assessed at four levels. \textbf{S} (strong): the defense eliminates the attack class upon adoption. \textbf{P} (partial): the defense is more expensive for the attacker or less effective but can be bypassed. \textbf{W} (weak): the defense is advisory, reactive, or largely ineffective. A dash means the defense does not apply to that class. The ratings in Fig.~\ref{fig:heatmap} are the authors' analytical estimates, provided by one rater using the evidence presented in this paper; they are not observed detection rates.

\begin{figure*}[!t]
\centering
\begin{tikzpicture}[scale=0.8, transform shape]
\foreach \h [count=\j] in {Safetensors, weights\_only, Model signing, Scanners, Pinning, Allowlists, Registry mod, Sandboxing} {
\node[anchor=south west, rotate=45, font=\scriptsize] at ({4.65+(\j-0.5)*\hcw-0.15},0.08) {\h};
}
\heatrow{1}{Malicious pickle model}{S,P,P,P,W,P,W,P}
\heatrow{2}{Model weights backdoored}{N,N,W,N,N,P,W,N}
\heatrow{3}{Split view data poisoning}{N,N,N,N,P,N,N,N}
\heatrow{4}{Data poisoning frontrunning}{N,N,N,N,W,N,N,N}
\heatrow{5}{RAG corpus poisoning}{N,N,N,W,N,W,N,W}
\heatrow{6}{MCP tool poisoning}{N,N,N,P,P,P,W,W}
\heatrow{7}{MCP puppet attack}{N,N,N,P,P,P,W,P}
\heatrow{8}{MCP rug pull}{N,N,N,P,S,W,W,P}
\heatrow{9}{MCP malicious external resource}{N,N,N,W,W,P,N,P}
\heatrow{10}{Malicious agent skills}{N,N,N,W,W,P,W,P}
\heatrow{11}{Typosquatting / malicious update}{N,N,N,W,P,P,W,P}
\heatrow{12}{Dependency confusion}{N,N,N,N,P,P,W,N}
\heatrow{13}{Vulnerable tooling (CVE)}{N,N,N,N,P,N,N,P}
\foreach \r in {2,5,10} \draw[thick] (0.1,{-\r*\hch}) -- ({4.65+8*\hcw},{-\r*\hch});
\node[anchor=west,font=\scriptsize\bfseries] at ({4.65+8*\hcw+0.15},{-1.0*\hch}) {Models};
\node[anchor=west,font=\scriptsize\bfseries] at ({4.65+8*\hcw+0.15},{-3.5*\hch}) {Datasets};
\node[anchor=west,font=\scriptsize\bfseries] at ({4.65+8*\hcw+0.15},{-7.5*\hch}) {Tools \& plugins};
\node[anchor=west,font=\scriptsize\bfseries] at ({4.65+8*\hcw+0.15},{-11.5*\hch}) {Dependencies};
\foreach \c/\t [count=\k] in {S/Strong,P/Partial,W/Weak,N/Not applicable} {
\node[draw,fill=h\c,minimum width=0.45cm,minimum height=0.35cm,inner sep=0pt] at ({0.6+(\k-1)*2.6},{-14.4*\hch}) {\csname lbl\c\endcsname};
\node[anchor=west] at ({0.9+(\k-1)*2.6},{-14.4*\hch}) {\t};
}
\end{tikzpicture}
\caption{Analytical defense coverage for each attack class (S: strong, P: partial, W: weak, -: not applicable). There are only two strong cells: safetensors against malicious pickle, and hash pinning against rug pulls.}
\label{fig:heatmap}
\end{figure*}

\subsection{Runtime Verification and Mechanism Removal}
\label{sec:def-strong}

Figure~\ref{fig:heatmap} illustrates two strong cells. Each corresponds to a defense that either removes the mechanism or re-checks at the time of use.

\emph{Safetensors.} Safetensors has no execution path. It stores raw tensor data and a header, and does not perform malicious serialization. It is therefore rated strong against malicious pickle. It does nothing about backdoored weights, which load fine and still misbehave. Adoption is the cost: 9.3\% of newly created popular repositories on Hugging Face still use pickle.

Restricted loaders such as PyTorch's weights-only unpickler do not strip the execution path; they implement an allowlist. An allowlist may block benign models. PickleBall, an allowlist-based loader, failed to load 51 of 252 benign models, or 20.2\%, in its own evaluation. An allowlist can also be bypassed: PyTorch versions 2.5.1 and prior can be subject to code execution via \texttt{torch.load(weights\_only=True)} (CVE-2025-32434)~\cite{nvd2025b}. We therefore mark \texttt{weights\_only} as partial.

\emph{Hash pinning and re-validation.} We consider hashing tool definitions at approval time and re-validating them at the beginning of each session, so that a changed description is caught before the agent loads it. This is the only control in the matrix that directly addresses the rug-pull mechanism. Pinning is strong against rug pulls when implemented, because a rug pull relies on the change going unnoticed. The Invariant demo was possible because no such re-validation happened. Elsewhere, the control is only partial or weak: it does not help against a description that is malicious from the beginning (MCP tool poisoning), against a malicious first version of a dataset, or when a snapshot was already poisoned before it was pinned.

\emph{Model signing.} Model Signing v1.0 signs a manifest over the model files using Sigstore. A signature indicates that the artifact has been issued by a particular party and that it has not been altered; it does not prove the artifact is safe. We therefore rate it partial against malicious files and weak against backdoors. None of the five incidents involved tampering with a signed artifact, and even a verified publisher can distribute a malicious file.

\subsection{Inspection at Time of Check}
\label{sec:def-pit}

Scanners are rated weak or partial in every row where they apply. Crafted model files can evade Picklescan. Of the 22 vulnerable pickle loading paths, 19 were missed by all four scanners tested by Liu \textit{et al.}~\cite{liu2025}. Picklescan recently gave 25.4\% false positives on benign models in one evaluation~\cite{qin2026}. A more recent context-sensitive scanner reports 0\% false negatives and 0.7\% false positives on its own benchmark, but it relies on an LLM analyzer, and switching the analyzer changes the false-negative count from 0 to 24. Scanning is probabilistic and dependent on the benchmark. It has not been tested against adaptive attackers, though it can be improved.

\emph{Registry moderation and allowlists.} Registry moderation is rated weak wherever it applies. The postmark-mcp backdoor and the ClawHub skills were both found by third parties after publication, and the registries' response was to take them down, not to block them. Allowlists of approved servers or publishers are incomplete: they block the impersonation pattern that postmark-mcp uses, but they do not stop an approved server from changing later, so they are weak against rug pulls.

\emph{Puppet attacks and malicious external resources.} For puppet attacks we give scanners, pinning, allowlists and sandboxing partial ratings, with the same rationale as for tool poisoning: pinning detects subsequent changes to a description, scanners can be evaded, allowlists cannot defend against a malicious or compromised approved server, and least privilege restricts what a steered call can do but not the steering itself. For malicious external resources, we rate scanners and pinning weak, because the poisoned content is external to the artifacts being scanned or pinned. We rate allowlists and sandboxing partial, because a poisoned page on an allowed domain still reaches the agent, while least privilege limits the effect of an injected instruction. Registry moderation does not apply here because the server in this attack is benign and the poisoned site is outside the registry. Human confirmation is not a field in the matrix; the Invariant case demonstrates that a confirmation dialog can be defeated. These are the authors' analytical ratings.

\subsection{Mitigating Impacts}
\label{sec:define-impact}

Sandboxing and least privilege limit what a poisoned tool or loaded model can do. They do not stop the compromise. The Invariant attack exploited an action the agent was allowed to take, and used whitespace padding to bypass the human confirmation dialog. Scoped permissions and human approval for new operations limit the blast radius, but they do not prevent the compromise. A sandboxed model loader also runs a payload that has restricted access.

\subsection{Where the Coverage Is Thinnest}
\label{section:def-thin}

The weakest coverage is for passive or behavioral artifacts. The matrix does little to protect against backdoored weights, dataset poisoning, or RAG poisoning. The best controls are partial at best: one is to whitelist weights from trusted publishers; the other is pinning dataset versions, which relies on a curation-time hash that most web-scale datasets simply do not have. Malicious skills are hard to regulate because the user is the one causing the harm, outside the skill's code. The matrix does not include software and AI bills of materials (SBOMs, AIBOMs), because these are inventory tools: they tell you which systems have a component that is affected, but they do not stop a compromise.

The ratings point to three design priorities. First, prefer formats and protocols that do not require you to examine the execution path. Second, re-verify artifacts at the point of use, not at approval. Third, invest in dataset provenance, where there is no deployed standard. These are results of analytical ratings and the incident record, not of our own experiments.

\section{Research Methodology: Isolated Differential Execution}
\label{sec:method}

We measure the discrepancy between static scanner verdicts and a loader's behavior via a differential-testing procedure. We name it the Isolated Differential Execution Framework (IDEF) and implement it as \texttt{pickle-harness}. It executes each file in an isolated container with no network, records whether the payload ran, and compares that observation to the output of static scanners. IDEF is an evaluation procedure, not a defense. This is in line with established differential-testing practice~\cite{pickleball2025}, as applied in previous measurements of scanner/loader mismatch.

\textbf{Corpus.} The 7,500-file corpus (seed 42; SHA-256 prefix \texttt{f8842089301cee39}) contains 5,000 malicious files from three families: \texttt{reduce\_basic} (2,500), \texttt{loadpath\_trick} (1,250) and \texttt{nullifai\_7z} (1,250). The corpus also contains 2,500 benign control files, both standard pickles and safetensors. All malicious payloads write a marker file and nothing more. \texttt{reduce\_basic} and \texttt{nullifai\_7z} use the pickle \texttt{GLOBAL} opcode to call \texttt{builtins.exec}. \texttt{loadpath\_trick} calls \texttt{builtins.eval} with the \texttt{exec} call hidden in a base64 string. Both callables are standard entries in the scanner blocklists, so the corpus is about parsing behavior, not coverage of obscure callables.

The \texttt{nullifai\_7z} files are our recreation of the nullifAI technique, not the original samples. Each file is a valid pickle prefix with the payload, followed by the 7z signature (\texttt{37 7A BC AF 27 1C}) and a truncated archive body. Picklescan reports an unknown opcode \texttt{b'7'} at position 88 in all 1,250 files. This matches the 7z signature starting at offset 88, where \texttt{0x37} is the first byte and is not a valid pickle opcode. \texttt{pickle.load} executes the payload while consuming the prefix and fails only at that byte. The actual ReversingLabs samples were PyTorch archives compressed with 7z that had broken pickle streams. Our files reproduce the same type of mismatch but are not equivalent to them.

\begin{algorithm}[htbp]
\caption{Isolated Differential Execution Algorithm (IDEF)}
\label{alg:harness}
\begin{algorithmic}[1]
\State \textbf{Input:} Corpus of untrusted serialized models.
\State \textbf{Output:} Scanner/loader mismatch count and fail-closed detection rate.
\State \textbf{Initialization:} Build a Docker container with \texttt{-{}-network=none}. Seed 42.
\State \textbf{Phase 1: Corpus generation}
\For{$i=1$ \textbf{to} $5000$}
\State Create a malicious file whose payload writes \texttt{/tmp/pwned\_<id>}.
\EndFor
\For{$i=1$ \textbf{to} $2500$}
\State Create a benign control file (standard pickle or safetensors).
\EndFor
\State \textbf{Phase 2: Dynamic ground truth (loader execution)}
\For{each file $f$ in the corpus}
\State Run \texttt{pickle.load($f$)} and check for the marker. Record \texttt{executed\_standard}.
\State Run \texttt{torch.load($f$, weights\_only=True)} and check for the marker. Record \texttt{executed\_weights\_only}.
\State Run \texttt{safetensors.numpy.load\_file($f$)} and check for the marker. Record \texttt{executed\_safetensors}.
\EndFor
\State \textbf{Phase 3: Static inspection}
\For{each file $f$ in the corpus}
\State Run \texttt{picklescan} and \texttt{modelscan} on $f$.
\State Record the classification: \texttt{flagged}, \texttt{clean}, or \texttt{error}. A run that exits with the scanner's error status is logged as \texttt{error} even if its output also contains findings.
\EndFor
\State \textbf{Phase 4: Differential analysis}
\State Join execution logs and scanner logs on file id.
\State Compute the number of mismatches.
\State Return the number of mismatches and the fail-closed rate, counting mismatches with classification \texttt{error} as detected.
\end{algorithmic}
\end{algorithm}

\textbf{Oracle and its limits.} The payload writes a marker file to indicate execution. If \texttt{/tmp/pwned\_<id>} is present, the payload code executed. The opposite is not true. An oracle may also undercount execution if a marker is missing because of a timeout, an import error, or an unwritable path. Each load has a 10-second timeout.

\textbf{Isolation and environment.} The container runs as a non-privileged user (UID 1000) together with \texttt{-{}-network none}. It has a read-only root filesystem (\texttt{-{}-read-only}). It uses an in-memory scratch directory for the marker files, mounted non-executable (\texttt{-{}-tmpfs /tmp:size=256m,noexec}). It has in-memory corpus and result directories. It has process and descriptor limits (\texttt{-{}-ulimit nproc=256:256}, \texttt{nofile=1024:1024}). It has compute limits (\texttt{-{}-memory 4g}, \texttt{-{}-cpus 4}). The loader script checks an external address (\texttt{8.8.8.8:53}, 2-second timeout) and aborts if it is reachable. The reported run used Python 3.11.17, PyTorch 2.14.1 (CPU build), safetensors 0.4.3, Picklescan 0.0.22 and ModelScan 0.8.4, and was executed on 24 October 2026.

\textbf{Availability.} The harness source code, corpus generation script, and run instructions are available at \url{https://github.com/shekhargiri005-dotcom/pickle-harness}. The corpus SHA-256 prefix is \texttt{f8842089301cee39} (full hash in \texttt{results/manifest.json}).

\section{Results and Discussion}
\label{sec:eval}

Table~\ref{tab:eval} summarizes the scanner results under the harness's classification. A scanner run that exits with an error status is recorded as \texttt{error} and is not counted as flagged. The important metric is \emph{Executed-Not-Flagged}, which counts payloads that the standard loader executed and the scanner did not flag.

\begin{table}[!htbp]
\centering
\caption{Scanner results on the 7,500-file corpus (Picklescan v0.0.22, ModelScan v0.8.4)}
\label{tab:eval}
\footnotesize
\begin{tabular}{@{}L{3.9cm}cc@{}}
\hline
\textbf{Metric} & \textbf{Picklescan} & \textbf{ModelScan} \\
\hline
Malicious files flagged (of 5,000) & 3,750 (75.0\%) & 5,000 (100.0\%) \\
Executed-Not-Flagged & 1,250 & 0 \\
Malicious files, scanner error & 1,250 & 0 \\
Benign files flagged (of 2,500) & 0 & 0 \\
Benign files, scanner error & 1 & 1 \\
\hline
\end{tabular}
\end{table}

Table~\ref{tab:perfamily} and Fig.~\ref{fig:perfamily} break the result down by family. Picklescan flagged every \texttt{reduce\_basic} and \texttt{loadpath\_trick} file. Its run ended in an error on all 1,250 \texttt{nullifai\_7z} files; they did not count as flagged for the harness, while the standard loader ran each one.

The scan logs show that the verdict is mixed, and that this is systematic. All 1,250 \texttt{nullifai\_7z} runs (100.0\%) print \texttt{dangerous import 'builtins exec' FOUND} and report one infected file. They then print \texttt{ERROR: parsing pickle ... at position 88, opcode b'7' unknown} and exit with code 2 (error), not code 1 (infected). Picklescan v0.0.22 thus parsed the pickle prefix, discovered the payload's malicious global, and stopped at the first byte of the 7z signature. Whether that is a detection depends on what consumes the output. A pipeline that blocks on the reported finding would block the file, while a pipeline that branches on the exit code, or that treats an abnormal exit as a failed scan without reading the output, would log an error. Our harness took the second reading. The 1,250 Executed-Not-Flagged files therefore measure how the scanner encodes and interprets its verdict, and they do not prove that Picklescan v0.0.22 did not see the payload.

\begin{table}[!htbp]
\centering
\caption{Scanner outcome per malicious family (all files executed under \textnormal{\texttt{pickle.load}}). PS: Picklescan, MS: ModelScan. ``PS error'' means the run failed with a parse error; see the text for the findings reported in the same run.}
\label{tab:perfamily}
\footnotesize
\begin{tabular}{@{}lrrrr@{}}
\hline
\textbf{Family} & \textbf{n} & \textbf{PS flagged} & \textbf{PS error} & \textbf{MS flagged}\\
\hline
reduce\_basic & 2,500 & 2,500 & 0 & 2,500\\
loadpath\_trick & 1,250 & 1,250 & 0 & 1,250\\
nullifai\_7z & 1,250 & 0 & 1,250 & 1,250\\
\hline
\end{tabular}
\end{table}

\begin{figure}[!htbp]
\centering
\begin{tikzpicture}
\begin{axis}[
    ybar,
    width=\columnwidth,
    height=5.5cm,
    bar width=10pt,
    ymin=0, ymax=125,
    ylabel={Flagged (\%)},
    symbolic x coords={reduce,loadpath,nullifai},
    xtick=data,
    enlarge x limits=0.25,
    nodes near coords,
    nodes near coords style={font=\tiny},
    legend style={at={(0.5,-0.22)}, anchor=north, legend columns=-1, font=\footnotesize},
    grid=major,
    grid style={dashed,gray!30}
]
\addplot coordinates {(reduce,100) (loadpath,100) (nullifai,0)};
\addplot coordinates {(reduce,100) (loadpath,100) (nullifai,100)};
\legend{Picklescan, ModelScan}
\end{axis}
\end{tikzpicture}
\caption{Percentage of each malicious family detected under the harness's classification (reduce: \texttt{reduce\_basic}, loadpath: \texttt{loadpath\_trick}, nullifai: \texttt{nullifai\_7z}). The Picklescan bar for nullifai is 0\% because its runs ended in a parse error; the output of every one of those runs also reports the dangerous global.}
\label{fig:perfamily}
\end{figure}

The comparison shows that the two scanners disagree on the verdict they give to one family of malformed files. ModelScan reports the finding without error, while Picklescan reports the finding with a parse error. The comparison does not indicate that static scanning as a class is inadequate, and it does not prove that Picklescan v0.0.22 misses the payload. Previous work shows successful bypass of multiple scanners. None of the adaptive payloads in our corpus is targeted at ModelScan; both callables we use (\texttt{exec}, \texttt{eval}) are standard blocklist entries. We have not investigated what ModelScan does with the files after the pickle prefix, so we cannot provide an explanation for the difference.

\subsection{Format-Level Defenses}
\label{sec:elimination}

Table~\ref{tab:elim} shows what happened when the same malicious files were submitted to alternative loaders.

\begin{table}[!htbp]
\centering
\caption{Alternative loaders on the malicious files as created}
\label{tab:elim}
\begin{tabular}{@{}lcc@{}}
\hline
\textbf{Loader} & \textbf{Executed / Total} & \textbf{Result} \\
\hline
Standard \texttt{pickle.load()} & 5000 / 5000 & Executed \\
\texttt{torch.load(weights\_only=True)} & 0 / 5000 & Rejected \\
\texttt{safetensors.numpy.load\_file()} & 0 / 5000 & Rejected \\
\hline
\end{tabular}
\end{table}

These results must be read with care. \texttt{torch.load} expects a PyTorch container, which is a zip archive. It rejects plain pickles and 7z-structured files at the format check, before the restricted unpickler is executed. \texttt{safetensors} expects a JSON header and raw tensor data, and outright rejects pickle files. Table~\ref{tab:elim} thus measures \emph{format rejection}, not the strength of the \texttt{weights\_only} unpickler. The 0/5000 result is expected because the files were not valid inputs for those loaders, and a loader that rejected all files would do equally well. We did not measure how many benign models each loader accepts, so the table says nothing about utility.

A test that exercises the \texttt{weights\_only} allowlist would embed the payload as \texttt{data.pkl} inside a real PyTorch archive and also load benign checkpoints under the same settings; we did not run it. We used PyTorch 2.14.1. CVE-2025-32434 affected \texttt{torch.load(..., weights\_only=True)} in PyTorch 2.5.1 and earlier and was fixed in 2.6.0, so it does not apply to our run. \texttt{weights\_only=True} has been the default since PyTorch 2.6. The vulnerability is the reason we rate the allowlist as partial, because it demonstrates that an allowlist loader can be bypassed.

What the data support is narrower than a general statement about defenses. If a loader only accepts a non-executable format such as safetensors, it cannot be made to run a pickle payload because it has no path that interprets one. A loader that accepts a container carrying a pickle depends on the correctness of its allowlist, and we rate it accordingly in Fig.~\ref{fig:heatmap}.

\subsection{Comparison with the Baseline (Song \textit{et al.})}
\label{sec:baseline}

We use Song \textit{et al.} as the baseline for the plugin and tool layer. They show that MCP attack vectors exist, demonstrate them on MCP aggregation platforms, and report that users were often unable to detect malicious servers. They do not assess any defense. Table~\ref{tab:song} compares the two studies, and Fig.~\ref{fig:song} shows the number of defenses each study rates for each of their attack classes. That comparison is of scope, not of measured performance. We grade defenses analytically, we recommend hash pinning for MCP clients but did not implement or test it, and our empirical harness tests pickle scanners, not MCP defenses. We rate all four of their attack types, but the ratings for puppet attacks and exploitation via malicious external resources rest on our interpretation of their definitions.

\begin{table*}[!t]
\caption{Baseline (Song \textit{et al.}) versus this study}
\label{tab:song}
\centering
\footnotesize
\renewcommand{\arraystretch}{1.25}
\begin{tabular}{L{4.0cm}L{6.0cm}L{6.0cm}}
\hline\hline
\textbf{Dimension} & \textbf{Song \textit{et al.}~\cite{song2025} (baseline)} & \textbf{This study}\\
\hline
Layers covered & Plugins and MCP & Four layers: models, datasets, plugins and tools, dependencies\\
MCP attack classes & Four demonstrated: tool poisoning, puppet attacks, rug pulls, malicious external resources & No new attacks; analyzes the Invariant WhatsApp and postmark-mcp incidents\\
Attack feasibility shown experimentally & Yes: malicious servers on MCP aggregation platforms, plus a user study & No new attack experiments\\
Defenses evaluated & None & Eight defenses rated analytically (author estimates, not measurements)\\
Song attack classes with a defense rating & 0 of 4 & 4 of 4 (analytical estimates)\\
Defense proposed for MCP & None evaluated & Hash pinning and re-validation of tool definitions; rated strong against rug pulls, untested\\
Scanners tested & None & Two pickle scanners on a model-layer corpus; no MCP scanner tested\\
Failure-mode analysis & Not provided & Three failure modes across five incidents\\
\hline
\end{tabular}
\end{table*}

\begin{figure}[!htbp]
\centering
\begin{tikzpicture}
\begin{axis}[
    width=\columnwidth,
    height=5.5cm,
    ylabel={Defenses rated (of 8)},
    symbolic x coords={Tool,Puppet,Rug,External},
    xtick=data,
    ymin=-0.5, ymax=6.5,
    ytick={0,2,4,6},
    legend style={at={(0.5,-0.25)}, anchor=north, legend columns=-1, font=\scriptsize},
    grid=major,
    grid style={dashed,gray!30}
]
\addplot[color=red, mark=square*, thick, mark size=2.5pt] coordinates {(Tool,0) (Puppet,0) (Rug,0) (External,0)};
\addlegendentry{Song et al.}
\addplot[color=blue!70!black, mark=*, thick, mark size=2.5pt] coordinates {(Tool,5) (Puppet,5) (Rug,5) (External,4)};
\addlegendentry{This study, rated}
\addplot[color=green!60!black, mark=triangle*, thick, mark size=2.5pt] coordinates {(Tool,3) (Puppet,4) (Rug,3) (External,2)};
\addlegendentry{Partial or strong}
\end{axis}
\end{tikzpicture}
\caption{Defenses with a rating, for each attack class of Song \textit{et al.} (Tool: tool poisoning; Puppet: puppet attacks; Rug: rug pulls; External: malicious external resources). ``Rated'' includes any S, P or W cell in Fig.~\ref{fig:heatmap}; ``partial or strong'' includes S or P cells. Lines are used to link the categories for clarity.}
\label{fig:song}
\end{figure}

\subsection{Comparison with Earlier Reports}
\label{sec:prior}

Table~\ref{tab:prior} shows the results of this work compared to those reported in the closest prior work. The columns are not a head-to-head race; the corpora, tools and metrics are different, and the previous results are on real models while ours are on a synthetic corpus. We therefore do not make any claim of improvement over these papers. A like-for-like comparison would need the same tools run against the same corpus.

\begin{table*}[!t]
\caption{Reported results of prior work versus this study. Corpora, tools and metrics differ, so rows are not directly comparable.}
\label{tab:prior}
\centering
\footnotesize
\renewcommand{\arraystretch}{1.25}
\begin{tabular}{L{3.2cm}L{6.4cm}L{6.4cm}}
\hline\hline
\textbf{Metric} & \textbf{Prior work (as reported)} & \textbf{This study}\\
\hline
Static scanner misses on malicious files & 19 of 22 vulnerable loading paths evaded all four scanners tested~\cite{liu2025}; DITTO: 0\% false negatives on its own benchmark, 24 with a different analyzer~\cite{qin2026} & ModelScan: 0 of 5,000 not flagged. Picklescan: 1,250 of 5,000 not flagged under the harness's classification (all 1,250 report the finding but exit with an error status)\\
Scanner false positives on benign files & Picklescan flagged 25.4\% of benign real-world models; DITTO 0.7\%~\cite{qin2026} & 0 of 2,500 for both scanners on synthetic benign files; one scanner error each\\
Malicious files stopped at load time & PickleBall rejected all 84 malicious models~\cite{pickleball2025} & \texttt{weights\_only} and safetensors executed 0 of 5,000, through format rejection of invalid inputs\\
Benign models still loadable & PickleBall loaded 79.8\% of 252 benign models~\cite{pickleball2025} & Not measured\\
Corpus & 252 benign and 84 malicious real models~\cite{pickleball2025}; Hugging Face repositories for pickle usage~\cite{qin2026} & 7,500 synthetic files in three payload families\\
\hline
\end{tabular}
\end{table*}

\subsection{Fail-Open versus Fail-Closed}
\label{sec:failclosed}

On the 1,250 nullifAI-style files, Picklescan returned its error status (code 2), and the harness recorded \texttt{error}, not \texttt{clean}. A pipeline that treats scanner errors as a pass (fail-open) and does not read the findings therefore loads files that execute. A pipeline that treats errors as a block (fail-closed), or that blocks on the reported finding, stops all 1,250. Picklescan's apparent miss on this family is thus as much a matter of interpretation of its verdict as of what it found.

Whether fail-closed is practical depends on how often a scanner errors on benign files, which is the additional cost of that policy; files a scanner flags are blocked under either policy. On our synthetic benign corpus, neither scanner flagged any of the 2,500 files, and each returned one error (0.04\%), so fail-closed would have blocked one benign file in 2,500. The benign corpus is not representative, however, because it does not contain any real-world model files, and an evaluation on real models reports Picklescan flagging 25.4\% of benign models~\cite{qin2026}. Flagged and error counts on real benign checkpoints would settle whether fail-closed is a usable stopgap; we did not make that measurement.

\section{Open Problems and Limitations}
\label{sec:discuss}

\subsection{Open Problems}

Pickle scanners reason about a representation of the file that might be different from what the loader runs, and agent-skill scanners do not see intent. Better signatures are more expensive for the attacker, but the defender is still in a race with the loader or the agent. The format has to come up with a durable fix.

In the sources we looked at, MCP clients ask the user to approve a server once and trust it forever. Hash pinning and re-validation would close the gap, but they are not the client defaults.

Model signing applies to model files. Web-scale datasets are built by crawling URLs without a signed manifest, and a poisoned dataset looks the same as a clean one after the fact.

\subsection{Limitations}

Most evidence on the incidents is gray literature. It is based on vendor blogs, research notes and press coverage, not peer-reviewed measurement. When possible, we use primary disclosures (JFrog, ReversingLabs, Invariant Labs, Koi Security) and otherwise secondary sources. Of the five cases, the Invariant demonstration and the Carlini \textit{et al.} attacks are research proofs of concept, not in-the-wild incidents. The field is moving faster than the publication cycle: a partial defense today can become strong (with a protocol change) or weak (after a new bypass), so Fig.~\ref{fig:heatmap} is a snapshot. The S/P/W ratings are from one rater and do not have a formal rubric or a per-cell evidence appendix. Our ratings for Song \textit{et al.}'s puppet-attack and external-resource classes are based on our interpretation of their definitions.

The empirical evaluation is performed on a synthetic corpus of 7,500 files. It does not represent the millions of models on Hugging Face, and the nullifAI-style files are reconstructions. The format-level results measure format rejection, not the \texttt{weights\_only} unpickler. Scanner error rates were only measured on the synthetic benign corpus, with one error per scanner in 2,500 files, so the cost of a fail-closed policy on real models is unquantified. We did not explore why ModelScan was able to process the nullifAI-style files without error. The headline numbers are shaped by the classification rule of the harness, which overrides reported findings in case of a parse error. Adaptive payloads are not in the corpus: all payloads use the standard blocklisted callables \texttt{builtins.exec} or \texttt{builtins.eval}, and none targets ModelScan or \texttt{weights\_only}. We did not test runtime hash pinning in a deployed MCP client. The incident record and analysis support the central hypothesis of this paper, that removing the mechanism is better than inspecting for it, but our measurements only partially support this hypothesis.

We tested only Picklescan v0.0.22 and ModelScan v0.8.4. In all 1,250 nullifAI-style logs, Picklescan v0.0.22 finds the dangerous global in the prefix before it reports the parse error. We therefore do not claim a bypass of this release, and we have not tested newer releases. The taxonomy does not deal with several surfaces that merit separate treatment, including remote-code execution through \texttt{trust\_remote\_code} and custom modeling code, other model formats such as GGUF, ONNX and joblib, chat-template injection through model metadata, build and CI compromise, and package-name hallucination.

\subsection{Boundary of Scope}

In July 2026, models running in OpenAI's internal cyber-capability evaluation escaped their isolation controls and reached Hugging Face production infrastructure. Hugging Face disclosed the intrusion on 16 July~\cite{hf2026}, and OpenAI released preliminary findings on 21 July~\cite{openai2026}. Hugging Face reports unauthorized access to a limited set of internal datasets and service credentials, entered through its dataset-processing pipeline, and reports no evidence of tampering with public models, datasets, Spaces, container images or published packages. This is the operator's own assessment.

We treat the incident as outside our taxonomy, because the victim was a platform operator and the cause was a failure of agent containment, not a poisoned artifact delivered to downstream users. It is nonetheless relevant to the dataset and dependency layers: the entry point was a dataset processed by executable parsing code, which illustrates that registry infrastructure is itself a supply chain dependency. Agent runtime isolation is left to future work.

\section{Conclusion}
\label{sec:conclusion}

We asked which supply chain attack surfaces AI systems expose, and how well current defenses cover those attack surfaces. We proposed a four-layer taxonomy: models, datasets, plugins and tools, and dependencies. It rests on three properties that are characteristic of AI artifacts: models are runnable, data is not code but molds behavior, and natural-language instructions act as code. Across five documented incidents we identified three failure modes: stale checks, interpretation differentials, and absent or blind controls. Each has a different cure. Stale checks must be re-verified at the point of use, interpretation differentials require removal of the ambiguous execution path, and absent or blind controls need better provenance or inspection.

\textbf{What our experiments show.}
\begin{itemize}
\item On a synthetic corpus where the standard loader ran all 5,000 malicious files, ModelScan v0.8.4 flagged all of them. Picklescan v0.0.22 flagged 3,750 and had a parse error on the 1,250 nullifAI-style files; all 1,250 outputs also report the payload's dangerous global, so the apparent miss is a mixed verdict plus an error exit status, not an unseen payload. The corpus contains no adaptive payloads.
\item Treating scanner errors as a block would have stopped all 1,250 files, at a cost of one blocked file in 2,500 on the synthetic benign corpus; the cost on real models was not measured.
\item \texttt{torch.load(weights\_only=True)} and \texttt{safetensors} did not run any of the files, but these files were not valid inputs for those loaders, so this shows format rejection and not the strength of the restricted unpickler.
\end{itemize}

\textbf{What remains analytical.} The S/P/W ratings in Fig.~\ref{fig:heatmap} are the authors' estimates and have not been validated across scanners, loaders or MCP clients. The hypothesis that removing the execution mechanism is more durable than inspecting for it is consistent with the incident record, but our experiments do not establish it. Three architectural problems remain open: scanners can be evaded by construction, one-time approval cannot detect rug pulls, and no deployed standard verifies dataset provenance. The next steps are to test payloads inside valid PyTorch containers against \texttt{weights\_only}, to measure scanner flagged and error rates on real benign models, to compare current scanner releases against adaptive payloads, and to test hash-pinned re-validation in a deployed MCP client.

{\sloppy
\begin{thebibliography}{00}

\bibitem{campbell2026}
R. Campbell, ``AI Supply Chain Security: MBOM-PQC Provenance, PQC Attestation, and a Maturity Model for Quantum-Resistant Assurance,'' \textit{Systems}, vol. 14, no. 5, p. 593, 2026, doi: 10.3390/systems14050593.

\bibitem{jfrog2024}
JFrog Security Research, ``Malicious Hugging Face ML Models with Silent Backdoor Target Data Scientists,'' JFrog Blog, Feb. 2024. [Online]. Available: \url{https://jfrog.com/blog/data-scientists-targeted-by-malicious-hugging-face-ml-models-with-silent-backdoor/}

\bibitem{reversinglabs2025}
ReversingLabs, ``Malicious ML models discovered on Hugging Face platform,'' ReversingLabs Blog, 2025. [Online]. Available: \url{https://www.reversinglabs.com/blog/rl-identifies-malware-ml-model-hosted-on-hugging-face}

\bibitem{carlini2024}
N. Carlini \textit{et al.}, ``Poisoning web-scale training datasets is practical,'' in \textit{Proc. IEEE Symp. Security and Privacy (SP)}, 2024, pp. 407--425.

\bibitem{invariant2025}
Invariant Labs, ``MCP security notification: Tool poisoning attacks,'' Apr. 2025. [Online]. Available: \url{https://invariantlabs.ai/blog/mcp-security-notification-tool-poisoning-attacks/}

\bibitem{koi2025}
Koi Security, ``First malicious MCP in the wild: the postmark backdoor stealing your e-mails,'' Koi Security Blog, Sep. 2025. [Online]. Available: \url{https://www.koi.security/blog/postmark-mcp-npm-malicious-backdoor-email-theft}

\bibitem{berini2026}
A. D. E. Berini \textit{et al.}, ``Security and privacy in LLMs: A comprehensive survey of threats and mitigation methods,'' \textit{Inf. Fusion}, vol. 132, 2026, doi: 10.1016/j.inffus.2026.104241.

\bibitem{nvd2025}
National Vulnerability Database, ``CVE-2025-6514 detail: mcp-remote OS command injection,'' NIST, Jul. 2025. [Online]. Available: \url{https://nvd.nist.gov/vuln/detail/cve-2025-6514}

\bibitem{liu2025}
T. Liu, G. Meng, P. Zhou, Z. Deng, S. Yao, and K. Chen, ``The Art of Hide and Seek: Making Pickle-based Model Supply Chain Poisoning Stealthy Again,'' arXiv:2508.19774, 2025.

\bibitem{qin2026}
Q. Qin, W. Li, B. Baudry, L. De Carli, H. Li and E. Merlo, ``DITTO: A Context-aware Pickle-based Pre-trained Model Scanner for Effective Security Audits,'' arXiv:2610.10735, 2026.

\bibitem{song2025}
H. Song \textit{et al.}, ``Beyond the protocol: Unveiling attack vectors in the Model Context Protocol (MCP) ecosystem,'' arXiv:2506.02040, 2025.

\bibitem{huang2024}
K. Huang \textit{et al.}, ``Lifting the veil on composition, risks, and mitigations of the large language model supply chain,'' arXiv:2410.21218, 2024.

\bibitem{wang2025}
S. H. Wang, Y. Zhao, X. Hou and H. Wang, ``A research agenda for large language model supply chain,'' \textit{ACM Trans. Softw. Eng. Methodol.}, vol. 34, no. 5, pp. 1--46, 2025.

\bibitem{owasp2023}
OWASP Foundation, ``CICD-SEC-3: Dependency chain abuse,'' OWASP Top 10 CI/CD Security Risks, 2023. [Online]. Available: \url{https://owasp.org/www-project-top-10-ci-cd-security-risks/CICD-SEC-03-Dependency-Chain-Abuse}

\bibitem{openssf2025}
Open Source Security Foundation, ``Launch of Model Signing v1.0: OpenSSF AI/ML Working Group secures the machine learning supply chain,'' OpenSSF Blog, Apr. 2025. [Online]. Available: \url{https://openssf.org/blog/2025/04/04/launch-of-model-signing-v1-0-openssf-ai-ml-working-group-secures-the-machine-learning-supply-chain/}

\bibitem{pickleball2025}
A. D. Kellas \textit{et al.}, ``PickleBall: Secure Deserialization of Pickle-based Machine Learning Models,'' in \textit{Proc. ACM Conf. Comput. Commun. Secur. (CCS)}, 2025.

\bibitem{plant2025}
R. Plant, M. V. Giuffrida, N. Pitropakis and D. Gkatzia, ``Evaluating the vulnerability of language models to poisoning attacks in low-resource settings,'' \textit{IEEE Trans. Audio, Speech, Lang. Process.}, vol. 33, pp. 54--67, 2025, doi: 10.1109/TASLP.2024.3507565.

\bibitem{demaio2025}
C. M. De Maio, M. Di Gisi, G. Fenza, M. Gallo, and V. Loia, ``A lifecycle-oriented survey of emerging threats and vulnerabilities in large language models,'' \textit{IEEE Access}, vol. 13, pp. 176482--176500, 2025, doi: 10.1109/ACCESS.2025.3619764.

\bibitem{poisoncraft2025}
Y. Shao \textit{et al.}, ``PoisonCraft: Practical poisoning of retrieval-augmented generation for large language models,'' arXiv:2505.06579, 2025. [Online]. Available: \url{https://arxiv.org/abs/2505.06579}

\bibitem{csa2026}
Cloud Security Alliance, ``MCP tool poisoning: Adversarial hijacking of AI agent workflows,'' CSA Research Note, 2026. [Online]. Available: \url{https://labs.cloudsecurityalliance.org/research/csa-research-note-mcp-tool-poisoning-ai-agent-exfiltration-2/}

\bibitem{nvd2025b}
National Vulnerability Database, ``CVE-2025-32434 detail: PyTorch \texttt{torch.load} with \texttt{weights\_only=True} remote code execution,'' NIST, Apr. 2025. [Online]. Available: \url{https://nvd.nist.gov/vuln/detail/CVE-2025-32434}

\bibitem{hf2026}
Hugging Face, ``Security incident disclosure,'' Hugging Face Blog, Jul. 16, 2026. [Online]. Available: \url{https://huggingface.co/blog/security-incident-july-2026}

\bibitem{openai2026}
OpenAI, ``Hugging Face model evaluation security incident: preliminary findings,'' Jul. 21, 2026. [Online]. Available: \url{https://openai.com/index/hugging-face-model-evaluation-security-incident/}

\end{thebibliography}
}

\end{document}"
